%=====================================================================
\chapter{Asymptotic Notation}\label{ch:asymptotics}
%=====================================================================
\locator{1}{Part I --- Analysing an Algorithm}

\begin{outcomes}
By the end of this chapter you will be able to:
\begin{tight}
  \item \textbf{State} the formal definitions of $O$, $\Omega$, $\Theta$, $o$
        and $\omega$, and \textbf{explain} what each one quantifies over
        \emph{(CLO-5)}.
  \item \textbf{Prove} that a given function belongs to a given asymptotic
        class by exhibiting the witnesses $c$ and $n_{0}$, and
        \textbf{disprove} membership by contradiction.
  \item \textbf{Use} the limit test, and \textbf{recognise} the case in which it
        does not apply.
  \item \textbf{Order} the standard functions by growth, and \textbf{apply}
        transitivity and the sum rule to simplify a bound.
  \item \textbf{Identify} the four common abuses of the notation, and
        \textbf{state} what each one conceals.
\end{tight}
\end{outcomes}

\begin{prereq}
\chref{role} and \chref{cases}. From Appendix~B you need the logarithm
identities and, for Section~3.5 only, L'H\^opital's rule.
\end{prereq}

\begin{keyterms}
asymptotic upper bound $\bullet$ asymptotic lower bound $\bullet$ tight bound
$\bullet$ witness $\bullet$ $O$ $\bullet$ $\Omega$ $\bullet$ $\Theta$ $\bullet$
$o$ $\bullet$ $\omega$ $\bullet$ transitivity $\bullet$ transpose symmetry
$\bullet$ sum rule $\bullet$ limit test $\bullet$ trichotomy
\end{keyterms}

\section{Three Bounds, Defined}

\chref{role} said we discard constants and keep the dominant term. That is a
practice, not a definition, and the practice can be abused. This chapter gives
it a definition precise enough to prove things with.

Throughout, $f$ and $g$ map the positive integers to the non-negative reals ---
running times, so never negative.

\begin{definitionbox}[$O$, $\Omega$ and $\Theta$]
\textbf{Upper bound.} $O(g(n))$ is the \emph{set} of functions
\[ O(g(n)) = \bigl\{\, f(n) \;:\;
   \exists\, c > 0,\ \exists\, n_{0} > 0,\
   \forall\, n \ge n_{0},\ 0 \le f(n) \le c\,g(n) \,\bigr\}. \]

\smallskip
\textbf{Lower bound.}
\[ \Omega(g(n)) = \bigl\{\, f(n) \;:\;
   \exists\, c > 0,\ \exists\, n_{0} > 0,\
   \forall\, n \ge n_{0},\ 0 \le c\,g(n) \le f(n) \,\bigr\}. \]

\smallskip
\textbf{Tight bound.}
\[ \Theta(g(n)) = \bigl\{\, f(n) \;:\;
   \exists\, c_{1}, c_{2} > 0,\ \exists\, n_{0} > 0,\
   \forall\, n \ge n_{0},\
   0 \le c_{1} g(n) \le f(n) \le c_{2} g(n) \,\bigr\}. \]

\smallskip
The constants $c$ (or $c_{1}, c_{2}$) and $n_{0}$ are the \term{witnesses}. To
prove a bound you produce them; to disprove one you show none can exist.
\end{definitionbox}

\begin{alertbox}[``$f(n) = O(g(n))$'' is an abuse of the equals sign, and you should know it]
$O(g)$ is a set, so the honest statement is $f \in O(g)$. Everyone writes
$f(n) = O(g(n))$ anyway, including this book, and the convention is harmless as
long as you remember that it is \textbf{one-directional}:

\begin{tight}
  \item $n = O(n^{2})$ is true.
  \item $O(n^{2}) = n$ is meaningless.
  \item From $f = O(n^{2})$ and $h = O(n^{2})$ you \emph{cannot} conclude
        $f = h$.
\end{tight}

\smallskip
Read the ``$=$'' as ``is'' --- as in ``Rawalpindi is a city'', which does not
make every city Rawalpindi.
\end{alertbox}

\section{Proving a Bound From the Definition}

The definitions are existential, so a proof is a construction: name $c$, name
$n_{0}$, verify the inequality.

\begin{theorem}[A quadratic is $O(n^{2})$]
$3n^{2} + 10n + 7 = O(n^{2})$.
\end{theorem}

\begin{proof}
We must produce $c$ and $n_{0}$ with $3n^{2} + 10n + 7 \le cn^{2}$ for all
$n \ge n_{0}$. For $n \ge 1$ we have $n \le n^{2}$ and $1 \le n^{2}$, so
\[ 3n^{2} + 10n + 7 \;\le\; 3n^{2} + 10n^{2} + 7n^{2} \;=\; 20n^{2}. \]
Taking $c = 20$ and $n_{0} = 1$ satisfies the definition.
\end{proof}

\begin{conceptbox}[The witnesses are not unique, and need not be tight]
That proof produced $c = 20$. A sharper argument gives a smaller one: for
$n \ge 10$, $10n + 7 \le n^{2}$, so $c = 4$ and $n_{0} = 10$ also work.

\smallskip
\textbf{Both proofs are equally valid}, because the definition asks only that
\emph{some} pair exist. Students waste effort hunting for the smallest $c$; the
definition does not reward it. Pick the crudest bound that closes the
inequality --- here, replacing every term by the largest one.
\end{conceptbox}

\begin{lemma}[Polynomials]\label{lem:poly}
Let $p(n) = \sum_{i=0}^{d} a_{i} n^{i}$ with $a_{d} > 0$. Then
$p(n) = \bigTh{n^{d}}$.
\end{lemma}

\begin{proof}
\textbf{Upper bound.} Let $A = \sum_{i=0}^{d} |a_{i}|$. For $n \ge 1$ every
$n^{i} \le n^{d}$, so $p(n) \le A n^{d}$, giving $c_{2} = A$.

\smallskip
\textbf{Lower bound.} Write $p(n) = a_{d}n^{d} + r(n)$ where
$r(n) = \sum_{i<d} a_{i}n^{i}$. By the same argument $|r(n)| \le B n^{d-1}$
with $B = \sum_{i<d}|a_{i}|$. Hence
\[ p(n) \;\ge\; a_{d}n^{d} - Bn^{d-1}
   \;=\; n^{d}\Bigl(a_{d} - \frac{B}{n}\Bigr)
   \;\ge\; \frac{a_{d}}{2}\,n^{d}
   \quad\text{whenever } n \ge \frac{2B}{a_{d}}. \]
So $c_{1} = a_{d}/2$ and $n_{0} = \max\{1,\, 2B/a_{d}\}$ complete the pair of
bounds, and $p(n) = \bigTh{n^{d}}$ by definition.
\end{proof}

\begin{corollary}[Only the degree survives]
The degree of a polynomial running time is the only thing about it that
survives asymptotic analysis. $0.001n^{3}$ grows faster than $1000n^{2}$.
\end{corollary}

Disproof is the other direction, and it is a proof by contradiction.

\begin{theorem}[A quadratic is not linear]\label{thm:notlin}
$n^{2} \ne O(n)$.
\end{theorem}

\begin{proof}
Suppose it were. Then there exist $c > 0$ and $n_{0} > 0$ such that
$n^{2} \le cn$ for all $n \ge n_{0}$. Since $n > 0$ we may divide by $n$:
\[ n \le c \qquad \text{for all } n \ge n_{0}. \]
But $c$ is a fixed constant, so taking $n = \max\{n_{0},\, \lceil c \rceil +
1\}$ gives $n \ge n_{0}$ and $n > c$ simultaneously --- contradicting the line
above. No such $c$ exists.
\end{proof}

\begin{theorem}[$\Theta$ is exactly $O$ and $\Omega$ together]\label{thm:thetaboth}
$f(n) = \bigTh{g(n)}$ if and only if $f(n) = \bigO{g(n)}$ and
$f(n) = \bigOm{g(n)}$.
\end{theorem}

\begin{proof}
($\Rightarrow$) Suppose $f = \bigTh{g}$ with witnesses $c_{1}, c_{2}, n_{0}$.
Then $f(n) \le c_{2}g(n)$ for $n \ge n_{0}$, which is the definition of
$O$ with $c = c_{2}$; and $c_{1}g(n) \le f(n)$ for $n \ge n_{0}$, which is the
definition of $\Omega$ with $c = c_{1}$.

\smallskip
($\Leftarrow$) Suppose $f = \bigO{g}$ with witnesses $c_{u}, n_{u}$ and
$f = \bigOm{g}$ with witnesses $c_{\ell}, n_{\ell}$. Put $n_{0} =
\max\{n_{u}, n_{\ell}\}$. For every $n \ge n_{0}$ both inequalities hold at
once, so
\[ c_{\ell}\,g(n) \;\le\; f(n) \;\le\; c_{u}\,g(n), \]
which is the definition of $\Theta$ with $c_{1} = c_{\ell}$ and
$c_{2} = c_{u}$.
\end{proof}

\begin{conceptbox}[Why that theorem earns its place]
It is the most-used result in the subject and it is two lines long, so it is
tempting to skip. Do not: it is what licenses the standard division of labour
in every later analysis.

\smallskip
To establish a $\Theta$ bound you prove an upper bound and a lower bound
\emph{separately}, by completely different arguments. \chref{lowerbound} proves
$\Omega(n\lg n)$ for comparison sorting by a decision-tree argument that has
nothing in common with \chref{mergesort}'s $\bigO{n \lg n}$ recurrence --- and
Theorem~\ref{thm:thetaboth} is the reason the two halves combine into
$\bigTh{n\lg n}$.
\end{conceptbox}

\section{The Strict Bounds: \texorpdfstring{$o$ and $\omega$}{o and omega}}

$O$ allows $f$ and $g$ to grow at the same rate. The strict forms forbid it,
and the change in the definition is the position of one quantifier.

\begin{definitionbox}[$o$ and $\omega$]
\[ o(g(n)) = \bigl\{\, f(n) :
   \forall\, c > 0,\ \exists\, n_{0} > 0,\
   \forall\, n \ge n_{0},\ 0 \le f(n) < c\,g(n) \,\bigr\} \]
\[ \omega(g(n)) = \bigl\{\, f(n) :
   \forall\, c > 0,\ \exists\, n_{0} > 0,\
   \forall\, n \ge n_{0},\ 0 \le c\,g(n) < f(n) \,\bigr\} \]

\smallskip
\textbf{The whole difference is $\exists c$ against $\forall c$.} For $O$ it is
enough that \emph{one} constant works. For $o$, \emph{every} constant must
work, however small --- so $f$ must become negligible relative to $g$, not
merely bounded by it.
\end{definitionbox}

\begin{worked}[Testing the Quantifier]
\textbf{Is $2n^{2} = o(n^{2})$?} Take $c = 1$. We would need $2n^{2} < n^{2}$
for all large $n$, i.e. $2 < 1$. So \textbf{no} --- and the reasoning shows why:
$o$ requires the ratio to vanish, and $2n^{2}/n^{2} = 2$ for every $n$.

\smallskip
\textbf{Is $2n = o(n^{2})$?} Given any $c > 0$, we need $2n < cn^{2}$, i.e.
$n > 2/c$. So $n_{0} = \lceil 2/c \rceil + 1$ works. Since we produced an
$n_{0}$ \emph{for every} $c$, \textbf{yes}.

\smallskip
\textbf{The pattern.} $f = o(g)$ exactly when $f = \bigO{g}$ but
$f \ne \bigTh{g}$ --- the bound holds and is \emph{not} tight.
\end{worked}

\section{Properties, and the Analogy With Inequality}

\begin{theorem}[Transitivity]\label{thm:trans}
If $f(n) = \bigO{g(n)}$ and $g(n) = \bigO{h(n)}$, then
$f(n) = \bigO{h(n)}$.
\end{theorem}

\begin{proof}
Take witnesses $c_{1}, n_{1}$ for the first and $c_{2}, n_{2}$ for the second,
and let $n_{0} = \max\{n_{1}, n_{2}\}$. For $n \ge n_{0}$,
\[ f(n) \;\le\; c_{1}\,g(n) \;\le\; c_{1}\bigl(c_{2}\,h(n)\bigr)
   \;=\; (c_{1}c_{2})\,h(n), \]
where the second inequality may be multiplied by $c_{1} > 0$ without reversing.
So $c = c_{1}c_{2}$ and $n_{0}$ are witnesses for $f = \bigO{h}$. The same
argument, with the inequalities reversed, gives transitivity for $\Omega$, and
Theorem~\ref{thm:thetaboth} then gives it for $\Theta$.
\end{proof}

\begin{lemma}[Sum rule]\label{lem:sum}
If $f_{1} = \bigO{g_{1}}$ and $f_{2} = \bigO{g_{2}}$, then
$f_{1} + f_{2} = \bigO{\max\{g_{1}, g_{2}\}}$.
\end{lemma}

\begin{proof}
With witnesses $c_{1}, n_{1}$ and $c_{2}, n_{2}$, put $n_{0} = \max\{n_{1},
n_{2}\}$ and $M(n) = \max\{g_{1}(n), g_{2}(n)\}$. For $n \ge n_{0}$,
\[ f_{1}(n) + f_{2}(n) \le c_{1}g_{1}(n) + c_{2}g_{2}(n)
   \le c_{1}M(n) + c_{2}M(n) = (c_{1}+c_{2})M(n). \qedhere \]
\end{proof}

Lemma~\ref{lem:sum} is why ``keep the dominant term'' is legitimate: a program
that does an $\bigO{n^{2}}$ pass and then an $\bigO{n}$ pass is $\bigO{n^{2}}$,
and we do not have to think about it again.

\begin{conceptbox}[The analogy with $\le$, and the one place it fails]
\rowcolors{2}{white}{lightbg}
\begin{longtable}{@{}L{2.4cm}L{1.8cm}L{10.1cm}@{}}
\toprule
\rowcolor{primary!15}
\textbf{Asymptotic} & \textbf{Reads as} & \textbf{Properties it really has} \\
\midrule
\endhead
$f = \bigO{g}$   & $f \le g$ & transitive; reflexive ($f = \bigO{f}$) \\
$f = \bigOm{g}$  & $f \ge g$ & transitive; reflexive \\
$f = \bigTh{g}$  & $f = g$   & transitive, reflexive \emph{and} symmetric ---
  an equivalence relation, which the other four are not \\
$f = o(g)$       & $f < g$   & transitive; \emph{not} reflexive \\
$f = \omega(g)$  & $f > g$   & transitive; not reflexive \\
\bottomrule
\end{longtable}

\smallskip
\textbf{Transpose symmetry} also holds: $f = \bigO{g}$ exactly when
$g = \bigOm{f}$, and $f = o(g)$ exactly when $g = \omega(f)$.

\smallskip
\textbf{But trichotomy fails.} For real numbers, exactly one of $a < b$,
$a = b$, $a > b$ holds. For functions there is a fourth possibility:
\emph{incomparable}. Let
\[ f(n) = n^{1 + \sin n}. \]
As $n$ runs over the integers $\sin n$ oscillates densely in $[-1, 1]$, so the
exponent oscillates in $[0, 2]$. Infinitely often $f(n)$ is near $n^{0} = 1$,
which kills any $\Omega(n)$ bound; infinitely often it is near $n^{2}$, which
kills any $O(n)$ bound. So $f$ is \textbf{neither $\bigO{n}$ nor
$\bigOm{n}$} --- the two functions simply cannot be ordered.

\smallskip
Pathological, and worth one paragraph: it is the reason the next section's limit
test needs a disclaimer, and the reason you may not argue ``$f$ is not
$\bigO{g}$, therefore $f$ is $\bigOm{g}$''.
\end{conceptbox}

\section{The Limit Test}

Constructing witnesses is reliable and slow. For most functions met in
practice, a limit settles the question at once.

\begin{theorem}[Limit test]\label{thm:limit}
Suppose $L = \displaystyle\lim_{n \to \infty} \frac{f(n)}{g(n)}$ exists, with
$L = \infty$ allowed. Then
\[ L = 0 \implies f = o(g); \qquad
   0 < L < \infty \implies f = \bigTh{g}; \qquad
   L = \infty \implies f = \omega(g). \]
\end{theorem}

\begin{proof}
We prove the middle case; the others are similar. Let $0 < L < \infty$. By the
definition of a limit, taking $\varepsilon = L/2$ there is an $n_{0}$ such that
for all $n \ge n_{0}$,
\[ \Bigl| \frac{f(n)}{g(n)} - L \Bigr| < \frac{L}{2}
   \quad\Longrightarrow\quad
   \frac{L}{2} \;<\; \frac{f(n)}{g(n)} \;<\; \frac{3L}{2}. \]
Multiplying through by $g(n) > 0$ gives
$\tfrac{L}{2}g(n) < f(n) < \tfrac{3L}{2}g(n)$, so $c_{1} = L/2$ and
$c_{2} = 3L/2$ witness $f = \bigTh{g}$.
\end{proof}

\begin{alertbox}[The test is one-directional]
Theorem~\ref{thm:limit} assumes the limit \emph{exists}. If it does not, the
test says nothing at all --- it does not say the functions are incomparable.

\smallskip
For $f(n) = n^{1+\sin n}$ and $g(n) = n$ the ratio $n^{\sin n}$ oscillates
forever and the limit does not exist; here the functions do turn out to be
incomparable. But $f(n) = n(2 + \sin n)$ also has no limiting ratio against
$g(n) = n$, and yet $f = \bigTh{n}$ plainly, with $c_{1} = 1$ and $c_{2} = 3$.
When the limit fails, go back to the witnesses.
\end{alertbox}

\begin{worked}[Logarithms Lose to Every Power]
\textbf{Claim:} $\lg n = o(n^{\varepsilon})$ for every constant
$\varepsilon > 0$ --- however small.

\smallskip
\textbf{Step 1.} Both $\lg n \to \infty$ and $n^{\varepsilon} \to \infty$, so
the ratio is of the form $\infty/\infty$ and L'H\^opital's rule applies.
Treat $n$ as a real variable $x$ and use $\lg x = \ln x / \ln 2$:
\[ \lim_{x\to\infty} \frac{\lg x}{x^{\varepsilon}}
   = \lim_{x\to\infty} \frac{(1/(x \ln 2))}{\varepsilon x^{\varepsilon - 1}}
   = \lim_{x\to\infty} \frac{1}{\varepsilon \ln 2 \cdot x^{\varepsilon}}
   = 0. \]

\smallskip
\textbf{Step 2.} The limit is $0$, so by Theorem~\ref{thm:limit},
$\lg n = o(n^{\varepsilon})$.

\smallskip
\textbf{Step 3 --- what it means.} Take $\varepsilon = 0.01$. The claim is that
$\lg n$ is eventually negligible beside $n^{0.01}$, a function so flat it looks
constant on any plot you can draw. ``Eventually'' is doing heavy work: at
$n = 10^{6}$, $\lg n \approx 20$ while $n^{0.01} \approx 1.15$, so the
inequality has not begun to hold. It holds from about $n = 10^{530}$.

\smallskip
\textbf{Step 4 --- the lesson.} Asymptotic statements are about limits, and a
limit carries no promise about $n = 10^{6}$. This is the same lesson as
\chref{role}'s crossover at $n = 65$, in a more extreme form.
\end{worked}

\begin{conceptbox}[The hierarchy worth memorising]
Each function below is $o$ of the one after it, for any constants
$\varepsilon > 0$, $k \ge 1$ and $a > 1$:
\[ 1 \;\prec\; \lg\lg n \;\prec\; \lg n \;\prec\; (\lg n)^{k}
    \;\prec\; n^{\varepsilon} \;\prec\; n \;\prec\; n \lg n
    \;\prec\; n^{2} \;\prec\; n^{k} \;\prec\; a^{n} \;\prec\; n! \]

\smallskip
At $n = 10^{6}$, taking one operation per nanosecond:

\rowcolors{2}{white}{lightbg}
\begin{longtable}{@{}L{3.0cm}L{4.4cm}L{6.9cm}@{}}
\toprule
\rowcolor{primary!15}
\textbf{$g(n)$} & \textbf{Operations} & \textbf{Time} \\
\midrule
\endhead
$\lg n$       & 20                  & 20 nanoseconds \\
$n$           & $10^{6}$            & 1 millisecond \\
$n \lg n$     & $2.0 \times 10^{7}$ & 20 milliseconds \\
$n^{2}$       & $10^{12}$           & 17 minutes \\
$n^{3}$       & $10^{18}$           & 32 years \\
$2^{n}$       & ---                 & no physical comparison exists \\
\bottomrule
\end{longtable}

\smallskip
The gap from $n\lg n$ to $n^{2}$ --- 20 milliseconds to 17 minutes --- is the
one most algorithm design is about.
\end{conceptbox}

\section{What \texorpdfstring{$\Theta$}{Theta} Asserts, Measured}

$f = \bigTh{g}$ says $f(n)/g(n)$ is eventually trapped between two positive
constants. That is a claim about a ratio, so it can be tested directly --- and
tested against \emph{rival} exponents, which is the sharper experiment.

\begin{worked}[Ruling Out the Exponents Either Side]
\texttt{code/ch03\_fit.cpp} times insertion sort and divides by three candidate
functions. Each column is normalised to $1.000$ at the first row, because the
constant itself is machine-specific and the question is only whether the column
\emph{stays put}. Median of nine runs.

\rowcolors{2}{white}{lightbg}
\begin{center}
\begin{tabular}{@{}rrrrr@{}}
\toprule
\rowcolor{primary!15}
$n$ & $T$ (ms) & $T/n^{1.9}$ & $T/n^{2.0}$ & $T/n^{2.1}$ \\
\midrule
  4{,}000 &    2.45 & 1.000 & 1.000 & 1.000 \\
  8{,}000 &   10.22 & 1.116 & 1.041 & 0.972 \\
 16{,}000 &   39.19 & 1.146 & 0.998 & 0.869 \\
 32{,}000 &  162.82 & 1.276 & 1.036 & 0.842 \\
 64{,}000 &  678.31 & 1.424 & 1.079 & 0.818 \\
128{,}000 & 2595.28 & 1.460 & 1.032 & 0.730 \\
\bottomrule
\end{tabular}
\end{center}

\smallskip
\textbf{Read the three columns as three hypotheses.}
\begin{tight}
  \item $T/n^{1.9}$ \emph{rises} by 46\%, without turning over. If
        $T = \bigTh{n^{1.9}}$ this column would be trapped between constants;
        it is climbing, so $n^{1.9}$ is too small.
  \item $T/n^{2.0}$ stays inside $[0.998,\, 1.079]$ --- a spread of 8\% across
        a 32-fold range of $n$. That is a ratio trapped between two
        constants, which is exactly the definition.
  \item $T/n^{2.1}$ \emph{falls} by 27\%, monotonically. So $n^{2.1}$ is too
        large: $T = o(n^{2.1})$, an upper bound that is not tight.
\end{tight}

\smallskip
\textbf{The experiment therefore does more than confirm the exponent} --- it
excludes the two exponents either side, which is what distinguishes evidence
for $\bigTh{n^{2}}$ from mere consistency with it.
\end{worked}

\begin{worked}[The Same Test for $n \lg n$, Which Is Harder]
$n \lg n$ sits only a whisker above $n$, so this test needs both a wider range
of $n$ and more careful timing.

\rowcolors{2}{white}{lightbg}
\begin{center}
\begin{tabular}{@{}rrrrr@{}}
\toprule
\rowcolor{primary!15}
$n$ & $T$ (ms) & $T/n$ & $T/(n \lg n)$ & $T/n^{1.2}$ \\
\midrule
  100{,}000 &   11.49 & 1.000 & 1.000 & 1.000 \\
  200{,}000 &   23.89 & 1.040 & 0.980 & 0.905 \\
  400{,}000 &   52.63 & 1.145 & 1.022 & 0.868 \\
  800{,}000 &  107.19 & 1.166 & 0.988 & 0.769 \\
1{,}600{,}000 &  237.89 & 1.294 & 1.043 & 0.743 \\
3{,}200{,}000 &  489.47 & 1.331 & 1.023 & 0.666 \\
6{,}400{,}000 & 1036.62 & 1.410 & 1.036 & 0.614 \\
\bottomrule
\end{tabular}
\end{center}

\smallskip
\textbf{$T/(n\lg n)$ stays within 4.3\% of 1} over a 64-fold range of $n$. The
rivals both fail: $T/n$ climbs 41\%, $T/n^{1.2}$ falls 39\%.

\smallskip
\textbf{Now the detail that makes this more than a curve fit.} If $T$ really is
$cn\lg n$, then $T/n = c\lg n$, so the $T/n$ column should itself grow like
$\lg n$. Over this range that predicts a rise of
\[ \frac{\lg(6{,}400{,}000)}{\lg(100{,}000)}
   = \frac{22.61}{16.61} = 1.361. \]
Measured: $\mathbf{1.410}$, within 4\% of the prediction. The drift in the
failing column is not noise --- it is the $\lg n$ factor, measured.

\smallskip
\textbf{And what it took to get there.} The first draft of this program
allocated a fresh vector inside the timed region. At these sizes the allocation
cost more than the sort it was supposed to be measuring, and the
$T/(n \lg n)$ column came out as $1.000, 1.134, 0.804, 1.342, 0.839$ --- noise,
from which nothing follows. The fix was to allocate once outside the loop,
refill with \texttt{std::copy}, and \emph{measure the refill separately and
subtract it}; the \texttt{copy} column in the program's output is that
measurement, and it is how you can tell the subtraction was small enough to
trust. Appendix~A collects the rest of these rules.
\end{worked}

\section{Four Abuses, and What Each Conceals}

\begin{alertbox}[1. Using $O$ when you mean $\Theta$]
``Insertion sort is $\bigO{n^{2}}$'' is true. So is ``insertion sort is
$\bigO{n^{100}}$'', and so is ``$\bigO{2^{n}}$'' --- an upper bound stays true
however loose it gets. If you know the bound is tight, say $\Theta$ and claim
the credit; if you do not, $O$ is honest and you should expect to be asked.

\smallskip
The reverse error is worse: reading somebody's $\bigO{\cdot}$ as if it were a
$\bigTh{\cdot}$, and concluding an algorithm is slower than it is.
\end{alertbox}

\begin{alertbox}[2. $\Omega$ as an excuse]
``My algorithm is $\bigOm{n}$'' is a boast about a \emph{lower} bound, which
for an algorithm is bad news, not good. Lower bounds are interesting about
\emph{problems} --- ``every comparison sort is $\bigOm{n\lg n}$'',
\chref{lowerbound} --- because they say no algorithm can do better. About your
own code, $\bigOm{\cdot}$ means ``at least this slow''.
\end{alertbox}

\begin{alertbox}[3. Dropping the case]
$\bigTh{\cdot}$ says nothing about which input you had in mind. Quicksort is
$\bigTh{n^{2}}$ in the worst case and $\bigTh{n \lg n}$ in expectation; both
are true and they differ by a factor of thousands. \textbf{Name the case every
time} --- the notation has no slot for it, so the sentence must.
\end{alertbox}

\begin{alertbox}[4. Forgetting that the constants went somewhere]
Asymptotic notation deliberately discards exactly the information that decides
which program to run at your actual $n$. From \texttt{ch03\_fit.cpp}:
$100n\lg n < n^{2}$ first holds at $n = \mathbf{997}$ --- not at $n = 10$.
Below a thousand elements the quadratic algorithm wins, and $\bigTh{\cdot}$
cannot tell you that.

\smallskip
The notation is a tool for comparing \emph{ideas}. Choosing an
\emph{implementation} needs the constants back, which means measuring --- the
other half of every chapter in this book.
\end{alertbox}

\begin{pitfall}
\begin{tight}
  \item \textbf{Writing $O(2n^{2})$ or $O(n^{2} + n)$.} Legal but illiterate:
        the point of the notation is that constants and lower terms are already
        gone. Write $\bigO{n^{2}}$.
  \item \textbf{Hunting for the smallest $c$.} The definition asks for
        existence. Any witness pair is a complete proof.
  \item \textbf{Concluding $f = \bigOm{g}$ from ``$f \ne \bigO{g}$''.}
        Trichotomy fails; $n^{1+\sin n}$ is the counterexample.
  \item \textbf{Using the limit test without checking the limit exists.}
        When it does not, the test is silent, not negative.
  \item \textbf{Comparing $\bigO{\cdot}$ bounds to rank two algorithms.} Both
        being $\bigO{n^{2}}$ tells you nothing about which is faster, nor even
        that either is quadratic.
  \item \textbf{Saying ``$\lg n$ is basically constant''.} It is $o(n^{\varepsilon})$
        for every $\varepsilon$, and it is also 20 at $n = 10^{6}$ and 30 at
        $n = 10^{9}$. Negligible in the limit is not the same as negligible.
  \item \textbf{Timing inside an allocation.} Not a notation error, but the
        reason the second worked example nearly reported nonsense.
\end{tight}
\end{pitfall}

\begin{lab}[Fitting a Curve, and Excluding the Rivals]
Reproduces both tables in Section~3.6 and the crossover figure. Expect
thirty minutes; part A takes about a minute of CPU at $n = 128{,}000$.

\begin{lstlisting}[language=C++]
// ch03_fit.cpp -- what Theta(g) actually asserts, tested.
// Build: cl /O2 /EHsc /std:c++17 ch03_fit.cpp        (see Appendix A)

// --- 1. Insertion sort against three candidate exponents -----------
// Only one of the three columns can stay constant, and which one
// does is the empirical content of "Theta(n^2)".
//
// Normalising to the FIRST row makes that row's noise everybody's
// noise, so the first row must not be a sub-millisecond timing:
// start at n = 4000 and take the median of nine runs.
double b19 = 0, b20 = 0, b21 = 0;
for (size_t n = 4000; n <= 128000; n *= 2) {
    auto base = random_input(n, 2024);
    std::vector<int> work(n);
    double t = time_ms([&]{
        std::copy(base.begin(), base.end(), work.begin());
        insertion_sort(work);
    }, 9);
    double x = static_cast<double>(n);
    double r19 = t / std::pow(x, 1.9);
    double r20 = t / std::pow(x, 2.0);
    double r21 = t / std::pow(x, 2.1);
    if (b19 == 0) { b19 = r19; b20 = r20; b21 = r21; }
    printf("%8zu %10.2f %10.3f %10.3f %10.3f\n",
           n, t, r19/b19, r20/b20, r21/b21);
}

// --- 2. Merge sort against n*lg(n) and its two rivals --------------
// n*lg(n) is only a whisker above n, so this is a far harder test
// than part 1, and it needs better hygiene to pass.
//
// The first draft allocated a fresh vector inside the timed region.
// At these sizes that allocation costs more than it measures -- the
// T/(n lg n) column came out as 1.000, 1.134, 0.804, 1.342, 0.839,
// which says nothing at all.
//
// The fix: allocate ONCE outside the timing and refill with
// std::copy, then measure the refill alone and subtract it.
for (size_t n = 100000; n <= 6400000; n *= 2) {
    auto base = random_input(n, 2024);
    std::vector<int> work(n), buf(n);
    double tc = time_ms([&]{
        std::copy(base.begin(), base.end(), work.begin());
    }, 11);
    double tt = time_ms([&]{
        std::copy(base.begin(), base.end(), work.begin());
        merge_sort(work, buf, 0, n);
    }, 11);
    double t = tt - tc;          // the sort, with the refill removed
    /* print tc so the reader can check the subtraction was small,
       then t/n, t/(n lg n) and t/n^1.2 normalised to the first row */
}

// --- 3. Where an asymptotically worse function is still smaller ----
// Theta is a statement about a limit, and the limit can be a long way
// off. Report the crossover rather than pretending there is none:
// 100 n lg n < n^2 requires 100 lg n < n.
for (size_t n = 2; n < 100000; ++n) {
    double x = static_cast<double>(n);
    if (100.0 * x * std::log2(x) < x * x) {
        printf("\n   crossover: n lg n first wins at n = %zu\n", n);
        break;
    }
}
\end{lstlisting}

\textbf{What to notice.} The \texttt{copy} column in part 2 is there to be
audited, not admired: if it ever approaches the \texttt{total} column, the
subtraction is doing too much work and the experiment should be redesigned
rather than reported.

\smallskip
\textbf{Then try to break it.} Add a $T/n^{2.05}$ column to part 1. It will
drift down, but only slightly --- and that is the honest limit of this method.
Timing can separate $n^{1.9}$ from $n^{2.1}$ over a 32-fold range; it cannot
separate $n^{2}$ from $n^{2.02}$, and no amount of care will make it. Say what
an experiment can resolve before you claim a result from it.
\end{lab}

\begin{checkpoint}
\begin{tightnum}
  \item Prove that $5n^{3} + 200n^{2} + 17 = \bigO{n^{3}}$ by exhibiting
        witnesses $c$ and $n_{0}$, and verify your inequality at $n = n_{0}$.
  \item Is $n^{2} = o(n^{2}\lg n)$? Justify your answer from the definition,
        not from the limit test.
  \item A student writes: ``my algorithm is $\bigO{n^{3}}$ and my colleague's is
        $\bigO{n^{2}}$, so mine is slower.'' Give the two distinct reasons this
        conclusion does not follow.
\end{tightnum}
\end{checkpoint}

\begin{chaptersummary}
\begin{tight}
  \item $O$, $\Omega$ and $\Theta$ are \emph{sets} of functions, defined by the
        existence of witnesses $c$ and $n_{0}$. ``$f = \bigO{g}$'' means
        $f \in \bigO{g}$ and may not be read backwards.
  \item To prove a bound, exhibit witnesses --- any witnesses. To disprove one,
        derive a contradiction, as for $n^{2} \ne \bigO{n}$.
  \item $f = \bigTh{g}$ iff $f = \bigO{g}$ and $f = \bigOm{g}$, which is what
        lets upper and lower bounds be proved by unrelated arguments and
        combined.
  \item $o$ and $\omega$ differ from $O$ and $\Omega$ only in having
        $\forall c$ where those have $\exists c$, which forces the ratio to
        vanish or diverge.
  \item All five are transitive; only $\Theta$ is an equivalence relation; and
        trichotomy fails --- $n^{1+\sin n}$ is comparable to $n$ in no
        direction at all.
  \item The limit test decides the question whenever the limit exists, and is
        silent, not negative, when it does not.
  \item Measured: $T/n^{2.0}$ for insertion sort held within 8\% across a
        32-fold range of $n$ while $T/n^{1.9}$ rose 46\% and $T/n^{2.1}$ fell
        27\%. For merge sort $T/(n\lg n)$ held within 4.3\% across a 64-fold
        range, and the drift in the failing $T/n$ column matched the predicted
        $\lg n$ growth to within 4\%.
  \item The notation discards precisely what decides which program to run:
        $100n\lg n$ does not beat $n^{2}$ until $n = 997$.
\end{tight}
\end{chaptersummary}

\begin{reviewq}
  \item \textbf{(MCQ)} $f(n) = \bigTh{g(n)}$ holds if and only if:
        (a)~$f = \bigO{g}$; (b)~$f = \bigOm{g}$; (c)~both of these;
        (d)~$f = o(g)$.
  \item \textbf{(MCQ)} Which is \emph{false}? (a)~$n = \bigO{n^{2}}$;
        (b)~$n^{2} = \bigO{n}$; (c)~$2^{n+1} = \bigO{2^{n}}$;
        (d)~$\lg n = o(n)$.
  \item \textbf{(Short)} State the formal definition of $\bigO{g(n)}$ and
        explain the role of each quantifier \emph{(CLO-5)}.
  \item \textbf{(Short)} Explain the difference between $\bigO{g}$ and $o(g)$
        in terms of the quantifier on $c$, and give a function that is in the
        first but not the second.
  \item \textbf{(Short)} Why is a lower bound a desirable result about a
        problem but an undesirable one about your own algorithm?
  \item \textbf{(Applied)} Prove $\bigTh{\cdot}$ is transitive, using
        Theorems~\ref{thm:thetaboth} and~\ref{thm:trans} rather than arguing
        from the definition directly.
  \item \textbf{(Applied)} You time an unknown algorithm and find that
        $T(n)/n^{2}$ falls steadily while $T(n)/(n\lg n)$ rises steadily across
        your whole range of $n$. What can you conclude, what can you not, and
        what would you measure next?
\end{reviewq}
